% Bibliography entries for chapters 01--07. Input inside the shared
% thebibliography environment; these keys are deliberately namespaced.
\bibitem{intro-boe} McLeay, Michael, Amar Radia and Ryland Thomas (14 March 2014).
\textit{Money creation in the modern economy}. Bank of England Quarterly
Bulletin, 54(1), 14--27. Institutional description of modern UK money creation,
including lending constraints; no current quantitative estimate used.
\url{https://www.bankofengland.co.uk/quarterly-bulletin/2014/q1/money-creation-in-the-modern-economy}.

\bibitem{intro-ets} European Commission, Directorate-General for Climate Action.
\textit{About the EU ETS}. Maintained institutional overview, ``How does the
EU ETS work?'' Accessed 19 September 2026. Used only for the cap-and-trade,
monitoring, reporting and allowance-surrender mechanism; no price, effectiveness
estimate or proposed reform treated as an adopted rule.
\url{https://climate.ec.europa.eu/areas-action/carbon-markets/about-eu-ets_en}.

\bibitem{intro-water} Food and Agriculture Organization of the United Nations.
\textit{Water scarcity}. Land, Soil and Water overview; physical, access and
infrastructure scarcity. Accessed 19 September 2026. No numerical scarcity
estimate used.
\url{https://www.fao.org/land-water/water/water-scarcity/en}.

\bibitem{intro-health} World Health Organization (6 May 2025).
\textit{World report on social determinants of health equity}.
ISBN 978-92-4-010758-8. Official publication overview used for the qualitative
relationship of social conditions and resource access to health inequities;
no effect size or individual medical recommendation inferred.
\url{https://www.who.int/publications/i/item/9789240107588}.

\bibitem{intro-ipcc} IPCC (2023). \textit{Climate Change 2023: Synthesis Report.
Summary for Policymakers}. In H. Lee and J. Romero (eds.). Section A.1.1 and
note 5: observed 2011--2020 global warming relative to 1850--1900. Section
B.1.1 and notes 28--31: conditional 2081--2100 warming under named SSP
scenarios. These are dated assessment values, not current-year measurements
or assigned probabilities of future social pathways.
\url{https://www.ipcc.ch/report/ar6/syr/summary-for-policymakers/}.

\bibitem{intro-ipbes} IPBES (2019). \textit{Summary for policymakers of the
global assessment report on biodiversity and ecosystem services}.
S. D\'iaz et al. (eds.). Section A5: the assessment estimate of roughly one
million threatened animal and plant species. See also IPBES Secretariat
(30 May 2019), \textit{A million threatened species? Thirteen questions and
answers}, which explains the estimate. Threatened is not equivalent to
already extinct, and the estimate is not a current census.
\url{https://www.ipbes.net/zh/node/35313}.

\bibitem{intro-resources} UNEP and International Resource Panel
(1 March 2024). \textit{Global Resources Outlook 2024}. Official report
overview: the conditional increase in global resource extraction by about
60 percent between 2020 and 2060 without urgent concerted action.
\url{https://www.unep.org/resources/Global-Resource-Outlook-2024}.

\bibitem{intro-homeostasis} McFarland, Jenny, Mary Pat Wenderoth,
Joel Michael, William Cliff, Ann Wright and Harold Modell (2016).
\textit{A conceptual framework for homeostasis: development and validation}.
Advances in Physiology Education, 40(2), 213--222.
\url{https://doi.org/10.1152/advan.00103.2015}.
This educational framework addresses organism-level homeostasis through
negative feedback; it is not a theory of economic or ecosystem control.
